\scanpage{43}
\mspage{4}

Picard - on peut dire que la classification ``au sens Picard'' est plus fine que la classification au sens des (Gr)-catégories. \deleted{\unsure{1}}

Associant à tout $X\handwritten{\in}\operatorname{Ob}C$ la symétrie de $X\otimes X$, interprété comme \deleted{étan} étant la tensorisation par un élément bien déterminé $s(X)$ de $\handwritten{\pi}_1(C)=\operatorname{Aut}(1_C)$, on trouve un\deleted{e} \deleted{application}\inserted{homomorphisme} canonique
\[
s(C)=s:\handwritten{\pi}_{\deleted{\unsure{1}}0}\handwritten{\longrightarrow}\,{}^{\unsure{\prime|\cdot}}{}_2(\handwritten{\pi}_{\deleted{\unsure{0}}1}),
\]
où $ {}_2(\handwritten{\pi}_{\deleted{\unsure{0}}1})$ désigne le sous-groupe des éléments $x$ de $\handwritten{\pi}_1$ tels que $2x=0$. Le résultat principal de la théorie de structure des catégories de Picard peut s'énoncer en disant que $C$ est déterminée(à équivalence \inserted{\handwritten{non canonique}} de catégories de Picard près) par la donnée des groupes commutatifs $\handwritten{\pi}_0$ et $\handwritten{\pi}_1$, et l'homomorphisme $s:\handwritten{\pi}_0\handwritten{\to}{}_2(\handwritten{\pi}_1)$. Dire que cet homomorphisme est nul signifie d'ailleurs que $C$ est une catégorie de Picard \uline{stricte}, et on retrouve le résultat de \handwritten{[~]} suivant lequel une catégorie de Picard stricte est déterminée (à équivalence près) par la connaissance de ses seuls invariants $\handwritten{\pi}_0$ et $\handwritten{\pi}_1$.

\uline{Remarques}. La théorie de loc. cit. nous permet, plus précisément, d'interpréter la 2-catégorie des catégories de Picard strictes en termes de la sous-catégorie pleine de la catégorie dérivée de (Ab) (cat. des groupes abéliens), formée des complexes $\handwritten{K_{\bullet}}$ tels que $H^{\deleted{\unsure{i}}}_{\handwritten{i}}(\handwritten{K_{\bullet}})=0$ pour $i\ne0,1$. Le résultat précédent, assez grossier, se réduit \inserted{alors} à la remarque qu'un objet de cette catégorie est déterminé \deleted{\unsure{à}}\handwritten{à} isomorphisme (non unique !) près par la connaissance de son $H^{\deleted{\unsure{1}}}_{\handwritten{0}}$ $(=\handwritten{\pi}_0)$ et $H^{\deleted{\unsure{0}}}_{\handwritten{1}}$ $(=\handwritten{\pi}_1)$. Il y aurait lieu de développer également pour les catégories de Picard non nécessairement strictes, et pour les (Gr)-catégories, une théorie de structure pour la 2-catégorie qu'ils forment. De plus, ces \deleted{résultats}\inserted{développements} encore dans les limbes, \deleted{\unsure{d}}tout comme ceux que \deleted{nous avons effectivement} nous venons d'esquisser, devraient être un jour \deleted{étendus en} étudiés dans le cadre des champs (en (Gr)-catégories, ou en catégories de Picard) tout comme dans\inserted{\handwritten{Deligne}} \handwritten{[~]}.

\deleted{\unsure{2}}3. \uline{Catégories de Picard enveloppantes}.

Soient $C$ une \deleted{\unsure{ca}} $\otimes$-catégorie AUC (par exemple une catégorie avec objet final et admettant des produits cartésiens de deux éléments), $S$ une partie de $\operatorname{Ob}C$. On se pose le problème universel d'envoyer (par un $\otimes$-foncteur AUC) $C$ dans une \inserted{$\otimes$-}catégorie \deleted{\unsure{$G$}}\inserted{AUC $G$}, de tel\inserted{\handwritten{le}} façon que pour tout $X\handwritten{\in}S$, $F(X)$ soit un objet \uline{inversible} de $G$. Le résultat principal ici est que ce problème admet toujours une solution, qu'on pourra appeler la $\otimes$-catégorie AUC ``de fractions'' par rapport à l'ensemble $S$, et noter \deleted{\unsure{e}}\handwritten{é}ventuellement $S^{\handwritten{-1}}C$. On trouve en fait que pour toute $G$, le foncteur naturel
% archive sheet number: 42
